\chapter{Координати, потік і поле}\label{ch:geo2}

\section{Циліндричні координати і осесиметрія}\label{sec:geo2-cyl}

Токамак описують у циліндричних координатах $(R,\tor,Z)$: $R$ --- відстань від вертикальної осі симетрії машини, $Z$ --- висота над екваторіальною площиною, $\tor$ --- тороїдальний кут навколо осі. Площину $\tor=\text{const}$ називають \emph{полоїдальним перерізом}; саме в ній малюють усі «перерізи плазми». Ідеальний токамак \emph{осесиметричний}: $\pa_\tor=0$ для всіх рівноважних величин. Тому вся геометрія рівноваги --- двовимірна, і її повністю задає одна функція двох змінних $\psi(R,Z)$. Порушення осесиметрії (скінченне число котушок TF, розд.~\ref{sec:geo6-ripple}; резонансні збурення, розд.~\ref{sec:geo3-thetastar}) розглядають як малі поправки до цієї картини.

\section{Полоїдальний потік і поле}\label{sec:geo2-psi}

\begin{outofcorpus}[подання осесиметричного поля]
Будь-яке осесиметричне поле з $\grad\cdot\Bvec=0$ можна записати як
\begin{equation}\label{eq:geo2-B}
\Bvec=\grad\psi\times\grad\tor+F\,\grad\tor,\qquad |\grad\tor|=\frac1R .
\end{equation}
Підставивши $\grad\tor=\vect{e}_\tor/R$ і $\vect{e}_R\times\vect{e}_\tor=\vect{e}_Z$, $\vect{e}_Z\times\vect{e}_\tor=-\vect{e}_R$, отримуємо
\begin{equation}\label{eq:geo2-comp}
\BR=-\frac1R\frac{\pa\psi}{\pa Z},\qquad \BZ=\frac1R\frac{\pa\psi}{\pa R},\qquad \Btor=\frac{F}{R},\qquad \Bpol=\frac{|\grad\psi|}{R}.
\end{equation}
Зміст $\psi$: потік $\BZ$ через горизонтальний диск радіуса $R$ на висоті $Z$ дорівнює $\int_0^R\BZ\,2\pi R'\,\dd R'=2\pi[\psi(R,Z)-\psi(0,Z)]$. Отже, $\psi$ --- полоїдальний потік \emph{на радіан} тороїдального кута (звідси одиниця Вб/рад). $F=R\Btor$ --- полоїдальна функція струму; у рівновазі $F=F(\psi)$~\cite[розд.~3]{Manual}.
\end{outofcorpus}

Найважливіший наслідок~\eqref{eq:geo2-B}: $\Bvec\cdot\grad\psi=0$. Силова лінія ніколи не перетинає поверхню $\psi=\text{const}$, тож кожна лінія рівня $\psi$ у перерізі, обернена навколо осі $Z$, дає \emph{магнітну поверхню} --- вкладений тор. Уся геометрія рівноваги --- це геометрія ліній рівня однієї функції.

У коді формули~\eqref{eq:geo2-comp} записано дослівно. Функцію $\psi$ задано на сітці EFIT і інтерпольовано бікубічним сплайном, похідні --- аналітичні похідні сплайна:

\code{viz/src/tokviz/equilibrium.py}{176}{189}{(поле з $\psi$ і $F$)}

\noindent Рядки 177--184 --- три компоненти~\eqref{eq:geo2-comp}; \texttt{dpsi\_dR}, \texttt{dpsi\_dZ} (рядки 167--171 файлу) --- похідні сплайна. Рядки 186--189: $\Btor=F/R$ зі \emph{сталим} $F$. Для DIII-D це наближення: у реальній плазмі $F$ залежить від $\psi$ (діа- і парамагнітний ефект струмів плазми), але в наборі даних $F(\psi)$ немає (п.~\ref{sec:geo2-F}). Рівновага Соловйова для JET (розд.~\ref{ch:geo6}) використовує самоузгоджену $F(\psi)$.

\section{Сітка EFIT і знак $\psi$}\label{sec:geo2-sign}

У наборі Fusion Equilibrium Challenge кожен кадр --- масив $\psi$ розміром $65\times65$ (рядки --- $Z$, стовпці --- $R$) з кроком $\Delta R=\SI{26,6}{мм}$, $\Delta Z=\SI{50,0}{мм}$ для DIII-D (значення з запуску \texttt{geom\_tour.py}). Сітка груба: на мале коло радіусом \SI{10}{см} навколо осі припадає лише кілька вузлів. Тому все, що потребує других похідних (кривина поверхонь біля осі, гесіан у X-точці, розд.~\ref{ch:geo5}), треба читати з поправкою на цю роздільність.

\emph{Знак $\psi$ не універсальний.} Формула~\eqref{eq:geo2-B} фіксує лише, що поле визначає $\psi$ з точністю до сталої; напрямки струму плазми і поля TF різні на різних машинах, і EFIT кожної машини має власну конвенцію. У DIII-D магнітна вісь --- \emph{мінімум} $\psi$, у MAST --- максимум. Проєкт зводить обидва випадки до одного множником:

\code{viz/src/tokviz/config.py}{39}{43}{(конвенція знака)}

\noindent Після множення на \texttt{AXIS\_SIGN} вісь завжди максимум, і всі пошукові алгоритми (вісь, LCFS) пишуть один раз. Той самий словник є в оцінювачі конкурсу (\texttt{fusion\_scoring/common.py}). Забутий знак --- не дрібниця: у проєкті збої вилучення LCFS при перенесенні DIII-D$\to$MAST, які спершу прийняли за фізичний сигнал (гіпотеза C8), виявилися саме артефактом знака $\psi$~\cite[розд.~4, R9]{Manual}.

\section{Магнітна вісь і нормований потік}\label{sec:geo2-axis}

Магнітна вісь --- точка, де $\grad\psi=0$ і $\psi$ досягає екстремуму (O-точка, розд.~\ref{ch:geo5}). Пошук іде у два кроки: грубий максимум $\texttt{axis\_sign}\cdot\psi$ на внутрішніх вузлах сітки, потім уточнення методом Нелдера--Міда на сплайні:

\code{viz/src/tokviz/equilibrium.py}{201}{225}{(пошук магнітної осі)}

\noindent Рядки 209--211: крайні два ряди вузлів виключено, щоб максимум на межі сітки (поле котушок) не «виграв» у справжньої осі. Рядки 214--218: цільова функція $-\texttt{axis\_sign}\cdot\psi$ на сплайні, зі штрафом за вихід із сітки. Рядки 220--225: Нелдер--Мід з допуском $10^{-7}$~м.

\begin{established}{VE1}
Знайдена вісь збігається з віссю EFIT до $\Delta R=\SI{0,0002}{мм}$, $\Delta Z=\SI{0,0001}{мм}$ при комірці $26{,}6\times\SI{50,0}{мм}$~\cite{RegViz}. Наш запуск дає $\Rax=\SI{1,67817}{м}$, $\Zax=\SI{0,02623}{м}$ проти EFIT \SI{1,67817}{м} і \SI{0,02623}{м}.
\end{established}

Далі потрібна межа. Замість того щоб шукати її самостійно, код бере $\psib$ як медіану $\psi$ уздовж контуру LCFS, який EFIT надає разом з рівновагою. Тоді нормування збігається з EFIT-івським:

\code{viz/src/tokviz/equilibrium.py}{227}{247}{(межа і нормований потік)}

\noindent Рядки 237--241 --- означення
\begin{equation}\label{eq:geo2-psin}
\psiN=\frac{\psi-\psiax}{\psib-\psiax},\qquad \psiN=0 \text{ на осі},\quad \psiN=1 \text{ на межі}.
\end{equation}
$\psiN$ не залежить ні від знака, ні від масштабу $\psi$. Саме тому нею позначають поверхні в усіх графіках і таблицях: «поверхня $\psiN=0{,}95$» означає одне й те саме для DIII-D і MAST. Для кадру 75: $\psiax=\SI{-0,2932}{Вб/рад}$, $\psib=\SI{0,0752}{Вб/рад}$; $\psi$ \emph{зростає} назовні, бо вісь DIII-D --- мінімум.

\begin{established}{VE2}
Контур LCFS з набору даних справді є лінією рівня $\psi$: розкид $\psi$ уздовж нього становить $8{,}2\cdot10^{-5}$, тобто 0,022\,\% діапазону потоку. Це виправдовує медіану як $\psib$~\cite{RegViz}.
\end{established}

\section{Тороїдальне поле, якого немає в даних}\label{sec:geo2-F}

$\psi$ дає лише полоїдальне поле. Для $\Btor$ потрібна $F$, а набір Fusion Equilibrium Challenge її не містить. Проєкт відновлює $F$ з того, що в наборі є, --- $\qnf$. Оскільки $q$ лінійне за $F$ (розд.~\ref{sec:geo3-q}), $F$ знаходять одним діленням; деталі --- в п.~\ref{sec:geo3-q}.

\begin{established}{VE3}
$F=\SI{3,23341}{Т.м}$ дає $\Btor(\Rax)=\SI{1,927}{Тл}$, тоді як реальне поле DIII-D становить 1,9--2,1~Тл. Жодної інформації про $\Bzero$ у розрахунок не подавали~\cite{RegViz}. На геометричному центрі $R_{\mathrm{geo}}=\SI{1,656}{м}$ маємо $\Btor=\SI{1,953}{Тл}$.
\end{established}

Корисне відчуття масштабу: на зовнішньому екваторі межі ($R\approx\SI{2,26}{м}$) $\Bpol=\SI{0,359}{Тл}$, а $\Btor=F/R\approx\SI{1,43}{Тл}$. Тобто силова лінія нахилена до тороїдального напрямку лише на $\arctan(0{,}359/1{,}43)\approx14^\circ$: вона робить кілька обертів навколо головної осі, поки обійде переріз один раз. Скільки саме --- каже $q$.

\begin{practicum}[поле і знак]
(1)~У \texttt{geom\_tour.py} замініть \texttt{machine="DIII-D"} на \texttt{"MAST"} для того самого кадру DIII-D. Що знайде \texttt{find\_axis()}? Чому? (2)~Візьміть MAST \texttt{mast\_shot\_28348.parquet}: чи збігається знайдена вісь з \texttt{efit\_r\_axis}, \texttt{efit\_z\_axis}? (3)~Порахуйте $\Bpol$ і $\Btor$ на внутрішньому екваторі ($R\approx\SI{1,05}{м}$) і кут нахилу лінії. Де лінія крутіша --- зовні чи всередині?
\end{practicum}
